% ===== PRÉAMBULE CANONIQUE — à coller VERBATIM en tête de CHAQUE .tex =====
\documentclass[11pt,a4paper]{article}
\usepackage[utf8]{inputenc}\usepackage[T1]{fontenc}\usepackage{lmodern}
\usepackage[french]{babel}
\usepackage{amsmath,amssymb,mathtools}\usepackage{icomma}
\usepackage[version=4]{mhchem}          % \ce{...} : équations & formules
\usepackage{siunitx}\sisetup{locale=FR} % \qty{}{}, \si{}, \ang{}
\usepackage{chemfig}                    % \chemfig{...} : structures développées
\usepackage{xcolor}\usepackage{colortbl}
\usepackage{tikz}\usepackage{pgfplots}\pgfplotsset{compat=1.18}
\usepackage[most]{tcolorbox}\usepackage{enumitem}\usepackage{array,booktabs}
\usepackage[a4paper,margin=2.2cm]{geometry}
\usetikzlibrary{arrows.meta,calc,angles,quotes,positioning,patterns,backgrounds,shapes.geometric}
\definecolor{primary}{RGB}{222,49,99}\definecolor{primarydark}{RGB}{150,22,55}
\definecolor{defcol}{RGB}{0,114,178}\definecolor{methcol}{RGB}{52,92,148}
\definecolor{excol}{RGB}{0,135,90}\definecolor{attncol}{RGB}{200,40,55}
\tcbset{boxbase/.style={enhanced,breakable,boxrule=0.9pt,arc=2.5pt,
  left=3mm,right=3mm,top=2.3mm,bottom=2.3mm,fonttitle=\bfseries,coltitle=white}}
% --- boîtes TUTORIEL (noms EXACTS, ne pas renommer) ---
\newtcolorbox{cle}[1][]{boxbase,colback=defcol!6,colframe=defcol,
  title={\textsc{À retenir}\ifx\relax#1\relax\relax\else\textmd{ : \itshape #1}\fi}}
\newtcolorbox{methode}[1][]{boxbase,colback=methcol!7,colframe=methcol,sharp corners,
  title={\textsc{Méthode}\ifx\relax#1\relax\relax\else\textmd{ --- \itshape #1}\fi}}
\newtcolorbox{exemplebox}[1][]{boxbase,colback=excol!5,colframe=excol,
  title={\textsc{Exemple}\ifx\relax#1\relax\relax\else\textmd{ : \itshape #1}\fi}}
\newtcolorbox{piege}[1][]{boxbase,colback=attncol!6,colframe=attncol,
  title={\textsc{Piège}\ifx\relax#1\relax\relax\else\textmd{ : \itshape #1}\fi}}
\newtcolorbox{remarque}[1][]{boxbase,colback=black!4,colframe=black!45,
  title={\textsc{Remarque}\ifx\relax#1\relax\relax\else\textmd{ : \itshape #1}\fi}}
% --- boîtes EXERCICE / CORRIGÉ (noms EXACTS, auto-comptées) ---
\newtcolorbox[auto counter]{exo}[1][]{boxbase,colback=primary!4,colframe=primary,
  title={\textsc{Exercice}~\thetcbcounter\ifx\relax#1\relax\relax\else\textmd{ --- \itshape #1}\fi}}
\newtcolorbox[auto counter]{corrige}[1][]{boxbase,colback=excol!4,colframe=excol,
  title={\textsc{Corrigé}~\thetcbcounter\ifx\relax#1\relax\relax\else\textmd{ --- \itshape #1}\fi}}
\newcommand{\R}{\mathbb{R}}\newcommand{\N}{\mathbb{N}}\newcommand{\Z}{\mathbb{Z}}\newcommand{\ds}{\displaystyle}
\begin{document}
\begin{corrige}[Fonction principale et fonction secondaire]
\begin{enumerate}[label=\alph*)]
  \item un $-\ce{OH}$ (\textbf{alcool}) et un $-\ce{COOH}$ (\textbf{acide carboxylique}).
  \item L'acide carboxylique est prioritaire ($-\ce{COOH} > -\ce{OH}$) : c'est la fonction principale (suffixe \textit{acide \dots-oïque}). L'alcool devient un préfixe \textbf{hydroxy-}.
  \item Chaîne à $3$~C, le carbone de l'acide est le n$^\circ$1, le $-\ce{OH}$ est en position 3 : \textbf{acide 3-hydroxypropanoïque}.
\end{enumerate}
\end{corrige}

\begin{corrige}[Numéroter pour le plus petit indice]
\begin{enumerate}[label=\alph*)]
  \item En partant de la droite ($\ce{CH3}$ voisin du $\ce{CHOH}$) : le $-\ce{OH}$ est en position \textbf{2}. En partant de la gauche : il est en position \textbf{4}.
  \item On retient le sens donnant le plus petit indice ($2 < 4$) : la molécule est le \textbf{pentan-2-ol}.
\end{enumerate}
\end{corrige}

\begin{corrige}[Nommer une molécule ramifiée]
La chaîne principale (la plus longue contenant le $-\ce{OH}$) compte $3$~carbones : c'est un \textbf{propan-1-ol}. Une ramification méthyle est portée par le carbone n$^\circ$2. Nom : \textbf{2-méthylpropan-1-ol}.
\end{corrige}

\begin{corrige}[Préfixes de substituants]
\begin{enumerate}[label=\alph*)]
  \item \ce{CH3-CO-CHCl-CH3} : cétone (fonction principale, sur C2) $+$ chlore sur C3 $\Rightarrow$ \textbf{3-chlorobutan-2-one}.
  \item \ce{OHC-CH2-CH(NO2)-CH3} : aldéhyde (C n$^\circ$1) $+$ groupe nitro sur C3 $\Rightarrow$ \textbf{3-nitrobutanal}.
\end{enumerate}
\end{corrige}

\begin{corrige}[Du nom à la structure]
\begin{enumerate}[label=\alph*)]
  \item La fonction principale est l'\textbf{acide carboxylique} ($-\ce{COOH}$), portée par le carbone n$^\circ$1. «\,propano\,» $\Rightarrow$ chaîne de $3$~carbones ; «\,2-méthyl\,» $\Rightarrow$ méthyle sur le carbone n$^\circ$2.
  \item Formule semi-développée : \ce{CH3-CH(CH3)-COOH}.
  \item Formule brute : \ce{C4H8O2}.
\end{enumerate}
\end{corrige}

\end{document}
